\documentclass[11pt]{article}
\usepackage[utf8]{inputenc}
\usepackage[T1]{fontenc}
\usepackage[spanish,es-nodecimaldot]{babel}
\usepackage[margin=1in]{geometry}
\usepackage{amsmath,amssymb,amsthm,booktabs,array,graphicx,pgfplots}
\usepackage[colorlinks=true,allcolors=blue!50!black]{hyperref}
\emergencystretch=1em
\pgfplotsset{compat=1.18}
\newcommand{\E}{\mathbb E}
\newcommand{\Var}{\operatorname{Var}}
\newcommand{\Cov}{\operatorname{Cov}}
\newcommand{\K}{\mathcal K}
\newtheorem{proposition}{Proposición}
\newtheorem{assumption}{Supuesto}
\title{Transparencia, sentimiento y fundamentos\\
\large Una corrección de la información parcial en Zhang y Zhang (2023)}
\author{David Julio Olano Silva Nevado\\IA y Modelamiento Económico, UP}
\date{2 de noviembre de 2026}
\begin{document}
\maketitle
\begin{abstract}
\noindent Estudio cuándo una señal pública más precisa aumenta la sensibilidad
esperada del precio de una acción a sus fundamentos, manteniendo fijo el
sentimiento. Audito la sección 6.2 de Zhang y Zhang (2023): una señal
determinista escalada revela todo el fundamento si su factor es conocido y
positivo, mientras que las varianzas impresas no recuperan los extremos
informativos. Propongo una señal gaussiana con ruido y derivo mediante Bayes
su componente público y residuo oculto. Desde utilidad CARA obtengo la
demanda óptima, el precio de equilibrio y una proposición de sensibilidad
creciente bajo condiciones explícitas. Código y simulaciones reproducen los
resultados. Lean certifica componentes algebraicos, con un alcance
probabilístico todavía parcial. La aplicación corresponde a acciones con
información desigual entre inversionistas, divulgación exógena y sin
inferencia adicional desde el precio.
\end{abstract}

\section{Pregunta y contribución}
\textquestiondown{}Cuándo una mayor precisión de la información pública hace
que el precio responda más a las noticias sobre el pago futuro de una acción?
Demuestro que la sensibilidad del precio esperado aumenta con la precisión
pública, bajo creencias normales, CARA, sentimiento fijo, oferta fija y sin
aprendizaje adicional desde precios. La conclusión trata de respuesta a
fundamentos, sin establecer aumentos universales del nivel ni del bienestar.

Conservo los cuatro grupos de Zhang y Zhang \cite{zhang2023} y reespecifico
su sección 6.2, ecuaciones (38)--(44), pp.~15--16 y proposición 5(ii).
La primera entrega postulaba las proporciones de varianza. Aquí las derivo
desde una señal ruidosa y Bayes, y obtengo la condición de primer orden
(CPO) desde la utilidad esperada CARA.

\section{Punto de partida y literatura}
El pago monetario futuro de una acción se representa de forma consistente por
\begin{equation}\label{eq:payoff}
 F=P_0+\theta+\epsilon.
\end{equation}
$P_0$ es una referencia conocida, $\theta$ la innovación fundamental y
$\epsilon$ una sorpresa final. Son cantidades monetarias por acción.
$\theta$ resume el efecto económico de la noticia sobre el pago futuro;
no es una cantidad de artículos ni una puntuación de sentimiento.

En el modelo básico de Zhang y Zhang las instituciones observan $\theta$
y los particulares no reciben información sobre su realización. Su ecuación
(15), p.~5, atribuye al particular racional varianza
$\sigma_\epsilon^2+\sigma_\theta^2$. En la extensión, la ecuación (39)
define $\kappa=\mu\theta$ y la ecuación (40) usa
$\sigma_\epsilon^2+\mu^2\sigma_\theta^2$. La proposición 5(ii) afirma
una mayor sensibilidad frente al modelo básico para cualquier $\mu$.

Si $\mu$ es conocido y positivo, $\theta=\kappa/\mu$: escalar el dato no
oculta su valor. En $\mu=0$, la varianza impresa elimina la incertidumbre
fundamental todavía presente; en $\mu=1$ la conserva aunque la noticia
ya se conoce. $\Var(\mu\theta)=\mu^2s$ es una identidad correcta para la
varianza marginal de la señal escalada. El problema es interpretarla como
incertidumbre restante después de observarla.

La auditoría de las fórmulas impresas admite un contraejemplo. Fijo
$\lambda_i=1/4$, $\gamma=R=1$, estados y sesgos sentimentales cero, y
$\sigma_\epsilon^2=s=\tau_1^2=\tau_2^2=1$. Estados cero no eliminan las
varianzas subjetivas. Entonces
\begin{equation}\label{eq:counter}
 K_0=\frac9{14}\simeq0{,}642857,\quad
 K^I(0)=\frac12,\quad K^I(0{,}1)=\frac{66943}{121303}\simeq0{,}551866<K_0.
\end{equation}
La comparación relevante es $K^I(0{,}1)$ frente al básico $K_0$.
$K^I(0)$ muestra además que la extensión no recupera su punto de partida.
El hallazgo concierne a las fórmulas e interpretación impresas; no identifica
si el origen fue editorial ni establece un efecto empírico de transparencia.

Di Maggio y Pagano \cite{dimaggio2018} modelan diferencias de procesamiento,
divulgación y decisiones del emisor, un mecanismo distinto del presente.
Grossman y Stiglitz \cite{grossman1980} muestran la importancia de información
costosa y aprendizaje desde precios; aquí restringir esa inferencia es un
supuesto explícito. El condicionamiento gaussiano es estándar. La contribución
es la reparación concreta y la derivación de sensibilidad de cuatro grupos.
La búsqueda delimitada de título, DOI, versiones y correcciones no sustenta
una afirmación de novedad universal. Los tres trabajos citados son artículos
publicados; el de Di Maggio y Pagano apareció en línea en 2017 y en volumen
en 2018.

\section{Aplicación y supuestos económicos}
La aplicación elegida es la negociación secundaria de una acción de una
empresa cotizada. La empresa es el emisor, pero no elige su divulgación
estratégicamente en el modelo. Un ejemplo es una noticia sobre demanda
futura: algunos inversionistas observan mejor su componente fundamental,
mientras el público recibe un reporte imperfecto. Observar exactamente
$\theta$ es una idealización de la ventaja institucional.

Información privada significa aquí que un dato pertenece al conjunto
informativo de unos inversionistas y no al de otros. Puede representar
análisis especializado u observación de indicadores. El modelo no identifica
el canal de filtraciones de directivos ni operaciones de políticos. Si solo
el emisor conoce algo y ningún inversionista lo observa, ambos grupos lo
desconocen: no corresponde asignarlo a una $\theta$ institucional conocida.

\begin{assumption}[Pago]\label{as:payoff}
$P_0$ es conocido, $\theta\sim N(0,s)$ con $s>0$,
$\epsilon\sim N(0,e)$ con $e=\sigma_\epsilon^2>0$; son independientes.
La oferta es una acción y $R=1+r_f>0$.
\end{assumption}
\begin{assumption}[Señal]\label{as:signal}
Para $0<\mu<1$, todos observan $y=\theta+\eta$, con
$\eta\sim N(0,t)$ independiente de $\theta,\epsilon$,
$t=s(1-\mu)/\mu>0$. Las instituciones observan además $\theta$.
Defino $\kappa=\E[\theta\mid y]$ y $\xi=\theta-\kappa$.
En $\mu=0$ defino $\kappa=0,\xi=\theta$; en $\mu=1$,
$\kappa=\theta,\xi=0$.
\end{assumption}
\begin{assumption}[Grupos y sentimiento]\label{as:sent}
Las masas $\lambda_i>0$ suman uno. Los grupos $1,2$ son institucionales y
$3,4$ particulares. Los grupos $2,4$ tienen estados $\rho_1,\rho_2$,
sesgos $\delta_1,\delta_2$ y varianzas subjetivas $\tau_1^2,\tau_2^2\geq0$.
Los estados son independientes de noticias y ruido; todos estos parámetros
permanecen fijos al comparar transparencia.
\end{assumption}
\begin{assumption}[Preferencias y creencias]\label{as:beliefs}
Cada creencia sobre el pago es normal, con media $m_i$ y varianza $v_i$.
$g_1=g_3=\gamma>0$, $g_2=\gamma e^{-b\rho_1}$,
$g_4=\gamma e^{-b\rho_2}$, $b>0$.
Los errores sentimentales subjetivos $\zeta_j\sim N(\delta_j,\tau_j^2)$
son independientes de los componentes informativos y la sorpresa final.
Describen creencias, no pagos adicionales reales. Las posiciones son reales
sin restricciones de venta corta ni endeudamiento.
\end{assumption}
\begin{assumption}[Información mantenida]\label{as:price}
Los particulares actualizan con $y$ y no extraen información adicional del
precio. Los agentes toman el precio como dado. No se modelan adquisición
costosa de información, elección del emisor ni reacción del sentimiento a
la divulgación.
\end{assumption}

La secuencia es: fijar parámetros, realizar $\theta$ y la señal, formar
creencias, negociar y vaciar el mercado, y realizar $\epsilon$. Los grupos
racionales carecen del sesgo sentimental especificado. El supuesto
\ref{as:price} limita una interpretación de expectativas racionales que
aprovechen toda la información del precio.

\section{Derivación de la información pública y oculta}
\subsection{Bayes y la distribución posterior}
La varianza de ruido $t$ mide imprecisión: un $t$ menor hace más útil el
reporte. Para un fundamento posible $x$ y un reporte observado $y$, Bayes da
\begin{align}
 f(x\mid y)&\propto f(y\mid x)f(x)\\
 &\propto \exp\left\{-\frac12
       \left[\frac{x^2}{s}+\frac{(y-x)^2}{t}\right]\right\}.
\end{align}
Expando y completo el cuadrado:
\begin{align}
 \frac{x^2}{s}+\frac{(y-x)^2}{t}
 &=\left(\frac1s+\frac1t\right)x^2-\frac{2y}{t}x+\frac{y^2}{t}\\
 &=\frac{s+t}{st}\left(x-\frac{s}{s+t}y\right)^2+\frac{y^2}{s+t}.
 \label{eq:posterior-square}
\end{align}
El último término no depende de $x$ y queda en la normalización. Reconocer
el núcleo normal entrega
\begin{equation}\label{eq:posterior}
 \theta\mid y\sim N\left(\frac{s}{s+t}y,\frac{st}{s+t}\right).
\end{equation}
Por eso defino
\begin{equation}\label{eq:mu}
 \mu=\frac{s}{s+t},\qquad \kappa=\mu y,\qquad \xi=\theta-\mu y.
\end{equation}
$\mu$ es adimensional. Ahora dividir $\kappa$ por $\mu$ recupera
$y=\theta+\eta$, que todavía contiene ruido. La señal es realmente parcial.

\begin{proposition}[Descomposición informativa]\label{prop:information}
Bajo los supuestos \ref{as:payoff}--\ref{as:signal}, $\theta=\kappa+\xi$.
Los componentes son gaussianos independientes, de media cero, y
\begin{equation}\label{eq:variances}
 \Var(\kappa)=\mu s,\quad \Var(\xi)=(1-\mu)s.
\end{equation}
Además,
\begin{equation}\label{eq:bridge}
 \E[\theta\mid\kappa]=\kappa,\quad
 \Var(\theta\mid\kappa)=(1-\mu)s,\quad
 \E[\kappa\mid\theta]=\mu\theta.
\end{equation}
\end{proposition}
\begin{proof}
Para el interior, $\Cov(\theta,\eta)=0$ por independencia y
$\Var(y)=s+t$. Aplicando $\Var(cX)=c^2\Var(X)$,
\begin{align}
 \Var(\kappa)&=\Var(\mu(\theta+\eta))=\mu^2(s+t)\\
 &=\frac{s^2}{(s+t)^2}(s+t)=\frac{s^2}{s+t}=\mu s.
\end{align}
Escribo $\xi=(1-\mu)\theta-\mu\eta$ y desarrollo todos los términos:
\begin{align}
 \Var(\xi)&=(1-\mu)^2s+\mu^2t-2\mu(1-\mu)\Cov(\theta,\eta)\\
 &=\frac{t^2s+s^2t}{(s+t)^2}
 =\frac{st(s+t)}{(s+t)^2}=\frac{st}{s+t}=(1-\mu)s.
\end{align}
La covarianza entre los componentes nuevos es
\begin{align}
 \Cov(\kappa,\xi)
 &=\Cov(\mu(\theta+\eta),(1-\mu)\theta-\mu\eta)\\
 &=\mu(1-\mu)s-\mu^2t
 =\mu[(1-\mu)s-\mu t]=0.
\end{align}
Son transformaciones lineales de normales independientes y, por tanto,
conjuntamente gaussianos. Para ese caso, covarianza cero implica
independencia; la implicación no es general fuera del caso gaussiano.
Así $\E[\xi\mid\kappa]=0$ y $\Var(\xi\mid\kappa)=\Var(\xi)$.
Sumar el componente conocido $\kappa$ prueba las dos primeras identidades
de \eqref{eq:bridge}. La última sigue de
$\E[\mu(\theta+\eta)\mid\theta]=\mu\theta$. Los extremos se verifican
directamente con las definiciones degeneradas del supuesto \ref{as:signal}.
$\mu=0$ representa ausencia de señal, no una normal con varianza infinita.
\end{proof}

La ley de varianza total permite interpretar la descomposición:
\begin{equation}
 s=\Var(\E[\theta\mid y])+\E[\Var(\theta\mid y)]
   =\mu s+(1-\mu)s.
\end{equation}
El primer término mide cómo varían los pronósticos públicos entre escenarios.
El segundo mide la incertidumbre que queda después del reporte. $\mu$
representa una fracción de varianza, no del valor realizado de $\theta$.

Por ejemplo, con $s=t=100$ se obtiene $\mu=1/2$. Si $y=8$ soles por acción,
$\kappa=4$: el pronóstico posterior es 4 y la varianza residual 50, con
desviación estándar $\sqrt{50}\simeq7{,}071$ soles por acción.
Si $P_0=100$ y $e=4$, el particular espera 104 con varianza 54.
Una institución que observa $\theta=10$ espera 110 con varianza 4.
La sorpresa $\epsilon$ todavía modifica el pago realizado.

\section{Varianzas y medias del pago percibido}
Una institución racional conoce $\theta$:
\begin{equation}
 m_1=P_0+\theta,\qquad v_1=\Var(\epsilon\mid\theta)=e.
\end{equation}
El particular racional conoce $\kappa$ y desconoce $\xi,\epsilon$:
\begin{align}
 m_3&=P_0+\kappa+\E[\xi\mid\kappa]+\E[\epsilon\mid\kappa]=P_0+\kappa,\\
 v_3&=\Var(\xi+\epsilon\mid\kappa)\\
 &=\Var(\xi\mid\kappa)+\Var(\epsilon\mid\kappa)
       +2\Cov(\xi,\epsilon\mid\kappa)=(1-\mu)s+e.
\end{align}
La covarianza condicional es cero por las independencias del modelo.
Añadir un error subjetivo independiente $\zeta_j$ desplaza la media en
$\delta_j$ y añade $\tau_j^2$ a la varianza. Con $d=(1-\mu)s$:
\begin{center}\small
\begin{tabular}{llll}
\toprule
Grupo & Media $m_i$ & Varianza $v_i$ & Aversión $g_i$\\
\midrule
1. Institución racional & $P_0+\theta$ & $e$ & $\gamma$\\
2. Institución sentimental & $P_0+\theta+\delta_1$ & $e+\tau_1^2$ & $\gamma e^{-b\rho_1}$\\
3. Particular racional & $P_0+\kappa$ & $e+d$ & $\gamma$\\
4. Particular sentimental & $P_0+\kappa+\delta_2$ & $e+d+\tau_2^2$ & $\gamma e^{-b\rho_2}$\\
\bottomrule
\end{tabular}
\end{center}
La empresa paga el mismo $F$ a todos. Estas son creencias distintas sobre
ese pago. Un desplazamiento constante $\delta_j$ no genera varianza:
la incertidumbre subjetiva adicional está en $\tau_j^2$.
Un estado $\rho_j=0$ implica $g_i=\gamma>0$; el agente sigue evitando riesgo.

\section{De utilidad CARA a la CPO}
\subsection{Riqueza, esperanza y varianza}
Con riqueza inicial $w_{0i}$ y posición $\omega_i$ al precio $P$,
\begin{align}
 W_i&=(w_{0i}-P\omega_i)R+\omega_i F\\
 &=Rw_{0i}+\omega_i(F-RP)=C_{0i}+\omega_i(F-RP).
\end{align}
Durante la elección, $P,\omega_i,g_i,m_i,v_i$ son constantes. Por linealidad,
\begin{equation}
 M_i=\E_i[W_i]=C_{0i}+\omega_i(m_i-RP).
\end{equation}
Sumar una constante no altera una varianza; escalarla multiplica por el
cuadrado del factor:
\begin{align}
 V_i=\Var_i(W_i)
 &=\Var_i(C_{0i}+\omega_i(F-RP))\\
 &=\omega_i^2\Var_i(F-RP)=\omega_i^2\Var_i(F)=\omega_i^2v_i.
\end{align}
Duplicar la posición cuadruplica la varianza de riqueza. CARA utiliza esas
varianzas para valorar riesgo; las varianzas provienen del pago y las creencias.

\subsection{Utilidad esperada y equivalente cierto}
La utilidad es $U_i(W)=-e^{-g_iW}$, con
$U_i'=g_i e^{-g_iW}>0$, $U_i''=-g_i^2e^{-g_iW}<0$ y
$-U_i''/U_i'=g_i$. La aversión absoluta es constante respecto de riqueza,
aunque difiere entre grupos o estados. La riqueza es normal bajo las
creencias del agente. Usando la identidad normal del apéndice,
\begin{align}
 \E_i[U_i(W_i)]
 &=-\exp(-g_iM_i+g_i^2V_i/2)\\
 &=-\exp[-g_i(M_i-g_iV_i/2)].
\end{align}
Como $z\mapsto-e^{-g_i z}$ es creciente, maximizar utilidad equivale a
maximizar el equivalente cierto
\begin{equation}\label{eq:ce}
 CE_i=C_{0i}+\omega_i(m_i-RP)-\frac{g_iv_i}{2}\omega_i^2.
\end{equation}
El equivalente cierto se mide en dinero y genera la misma elección.
No es una igualdad numérica con la utilidad esperada.

\begin{proposition}[Demanda óptima]\label{prop:demand}
Bajo los supuestos \ref{as:payoff}--\ref{as:price}, cada grupo tiene una
única posición óptima
\begin{equation}\label{eq:demand}
 \omega_i^*=\frac{m_i-RP}{g_iv_i}.
\end{equation}
\end{proposition}
\begin{proof}
En \eqref{eq:ce}, la derivada del ingreso fijo es cero; la del término
lineal es $m_i-RP$; la de la penalización es
$-(g_iv_i/2)2\omega_i=-g_iv_i\omega_i$. Por tanto,
\begin{align}
 CE_i'&=m_i-RP-g_iv_i\omega_i=0,\\
 g_iv_i\omega_i&=m_i-RP,\qquad
 \omega_i^*=(m_i-RP)/(g_iv_i).
\end{align}
Dividir es válido porque $g_i,v_i>0$. La segunda derivada es
$CE_i''=-g_iv_i<0$. Además,
\begin{equation}\label{eq:square}
 CE_i(\omega_i)=CE_i(\omega_i^*)-
 \frac{g_iv_i}{2}(\omega_i-\omega_i^*)^2.
\end{equation}
El término restado es no negativo y solo se anula en el óptimo:
hay máximo global único.
\end{proof}
El numerador es pago esperado menos costo financiado. El denominador
penaliza exposición por aversión y varianza. Para el particular racional,
\begin{equation}
 \omega_3^*=\frac{P_0+\kappa-RP}{\gamma[e+(1-\mu)s]}.
\end{equation}
La información cambia el pronóstico y la incertidumbre. Una posición
negativa representa venta corta permitida por el modelo.

\section{Precio de equilibrio y resultado propio}
Defino pesos $h_i=\lambda_i/(g_iv_i)>0$, $H=\sum_i h_i$,
$a=h_1+h_2$ y $B=h_3+h_4$. Reúnen masa y respuesta al riesgo, y no tienen
que sumar uno.
\begin{proposition}[Precio y sensibilidad esperada]\label{prop:price}
Bajo los supuestos \ref{as:payoff}--\ref{as:price}, el precio es único:
\begin{equation}\label{eq:price}
 P=\frac{HP_0+a\theta+B\kappa+h_2\delta_1+h_4\delta_2-1}{RH}.
\end{equation}
Con sentimiento fijo,
\begin{equation}\label{eq:K}
 \bar P=\E[P\mid\theta,\rho_1,\rho_2]=C(\mu)+\K(\mu)\theta,\quad
 \K(\mu)=\frac{a+\mu B}{R(a+B)},
\end{equation}
con $C(\mu)=[HP_0+h_2\delta_1+h_4\delta_2-1]/(RH)$.
\end{proposition}
\begin{proof}
La oferta de una acción exige
\begin{align}
 1&=\sum_i\lambda_i\omega_i^*
   =\sum_i h_i(m_i-RP)=\sum_i h_i m_i-RPH.
\end{align}
Despejar da $P=(\sum_i h_i m_i-1)/(RH)$, único porque $R,H>0$.
Sustituir las cuatro medias da \eqref{eq:price}. Los pesos no dependen de
las realizaciones de noticias. La independencia del sentimiento y la
proposición \ref{prop:information} dan
$\E[\kappa\mid\theta,\rho_1,\rho_2]=\mu\theta$, probando \eqref{eq:K}.
\end{proof}
$\K$ mide la respuesta del precio esperado al fundamento. Mantener fija
la señal realizada y derivar \eqref{eq:price} respecto de $\theta$ daría
$a/(RH)$, otra comparación. Si $\K=0{,}6$, una noticia adicional de una
unidad monetaria cambia el precio esperado en $0{,}6$ unidades.

\begin{proposition}[Mayor sensibilidad con sentimiento fijo]\label{prop:main}
Bajo los supuestos \ref{as:payoff}--\ref{as:price}, si
$0\leq x<y\leq1$, entonces $\K(x)<\K(y)$.
Además $0<\K(\mu)\leq1/R$ y $\K(1)=1/R$.
\end{proposition}
\begin{proof}
$a>0$ es constante porque las instituciones conservan varianzas y aversiones.
Sean $n_3=\lambda_3/\gamma>0$, $n_4=\lambda_4/g_4>0$,
$e_3=e>0$, $e_4=e+\tau_2^2>0$. Entonces
\begin{equation}
 d_j=e_j+(1-\mu)s>0,\quad
 B=\sum_{j=3}^4 n_j/d_j,\quad
 B'=\sum_{j=3}^4 n_js/d_j^2.
\end{equation}
La regla del cociente permite desarrollar
\begin{align}
 \K'&=\frac{(B+\mu B')(a+B)-(a+\mu B)B'}{R(a+B)^2}\\
 &=\frac{B^2+a[B-(1-\mu)B']}{R(a+B)^2}.
\end{align}
En cada sumando,
\begin{equation}
 \frac{n_j}{d_j}-(1-\mu)\frac{n_js}{d_j^2}
 =\frac{n_j[d_j-(1-\mu)s]}{d_j^2}
 =\frac{n_je_j}{d_j^2}>0.
\end{equation}
Por tanto
\begin{equation}\label{eq:derivative}
 \boxed{\K'=\frac{B^2+a\sum_{j=3}^4 n_je_j/d_j^2}{R(a+B)^2}>0.}
\end{equation}
La función es continua en el intervalo cerrado y diferenciable dentro:
el incremento es estricto. $a+\mu B\leq a+B$ prueba la cota superior,
la positividad de parámetros la inferior, y sustituir $\mu=1$ da $1/R$.
El apéndice ofrece una identidad de diferencias finitas.
\end{proof}
Con $\mu=0$ se recupera el coeficiente básico, porque los particulares
conservan toda la incertidumbre $s$. Con $\mu=1$ aún existe riesgo final
$e>0$. La prima por riesgo y los sesgos siguen afectando el nivel del precio.

\section{Verificación reproducible}
El punto de entrada es \texttt{python code/verify.py}, sin dependencias
externas. Comprueba el contraejemplo, identidades de posterior y momentos,
CPO, máximo y vaciado de mercado, y genera las cifras y datos de gráficos.
El barrido racional compara 5\,760 pares de transparencias y 60 mercados.
Una simulación gaussiana con semilla fija contrasta momentos y covarianzas.
También contrasta la expresión CARA con integración numérica determinista.
Las simulaciones verifican implementación; la prueba general es analítica.

\begin{center}
\begin{tabular}{lrrr}
\toprule
Escenario & $\mu=0$ & $\mu=0{,}1$ & $\mu=1$\\
\midrule
Básico, referencia constante & $0{,}642857$ & $0{,}642857$ & $0{,}642857$\\
Fórmula publicada & $0{,}500000$ & $0{,}551866$ & $1$\\
Corrección propuesta & $0{,}642857$ & $0{,}669346$ & $1$\\
\bottomrule
\end{tabular}
\end{center}
Los parámetros son los de \eqref{eq:counter}. La primera fila es una
referencia constante; no representa cambios informativos del básico.
El valor corregido en $0{,}1$ es exactamente $583/871$.
\begin{figure}[ht]
\centering
% Generated from checked data. No screenshots of published figures.
\begin{tikzpicture}
\begin{axis}[
 width=10.6cm,height=5.7cm,
 xmin=0,xmax=1,ymin=0.45,ymax=1.03,
 xlabel={Transparencia $\mu$},ylabel={Sensibilidad informativa $\mathcal K$},
 tick label style={font=\small},label style={font=\small},
 legend style={font=\scriptsize,draw=none,at={(0.02,0.98)},anchor=north west},
 grid=major,grid style={gray!15},
]
\addplot[blue!70!black,very thick] table[x=mu,y=corrected,col sep=comma]{../code/output/information.csv};
\addlegendentry{Corrección: precio esperado}
\addplot[orange!85!black,thick] table[x=mu,y=printed,col sep=comma]{../code/output/information.csv};
\addlegendentry{Fórmula publicada}
\addplot[black,dashed] table[x=mu,y=baseline,col sep=comma]{../code/output/information.csv};
\addlegendentry{Modelo básico}
\end{axis}
\end{tikzpicture}

\caption{Sensibilidad esperada: la corrección recupera el básico y aumenta
hasta $1/R$. Datos del código, sin estimación empírica.}
\end{figure}
Los informes \texttt{code/output/verification.json} y
\texttt{signal-verification.json} detallan recuentos y tolerancias.
El código exporta momentos y posterior para poder reproducir los ejemplos.

\section{Alcance, emisores y límites}
El modelo caracteriza un mecanismo bajo condiciones explícitas. No es una
estimación de una empresa particular ni una afirmación válida para todo
mercado. La misma acción tiene un pago común; un grupo conoce el fundamento
y otro recibe un reporte imperfecto. Aplicarlo a bonos soberanos requeriría
incorporar pago contractual, incumplimiento y sus determinantes.

Deben distinguirse tres casos. Información conocida solo por el emisor:
ambos grupos de inversionistas la desconocen. Información privada de algunos
inversionistas por análisis u observación: esta es la interpretación elegida,
idealizada como conocimiento exacto. Información transmitida desde un
directivo a un agente: requiere especificar transmisión, incentivos y
selección del receptor, elementos ausentes aquí. La palabra privada no
identifica por sí sola el origen de esa ventaja.

Tampoco se demuestra una ganancia extraordinaria de instituciones. La
demanda depende del precio y del riesgo final; una mejor previsión puede
modificar la posición sin garantizar una rentabilidad superior en cada
realización. Comparar rendimientos requiere un referente, costos y horizonte.
El objeto demostrado es sensibilidad $\K$, no un alfa de inversión.

El supuesto de no inferir del precio es material. Si los particulares
conocieran todos los pesos, sesgos y la oferta fija, podrían despejar
\begin{equation}
 \theta=\frac{RHP-HP_0-B\kappa-h_2\delta_1-h_4\delta_2+1}{a}.
\end{equation}
Por eso no presento esta solución como un equilibrio de expectativas
racionales con inferencia completa. Una extensión con oferta aleatoria o
señales institucionales imperfectas exigiría resolver creencias y precio
conjuntamente.

Si el sentimiento responde a la transparencia, cambian $a$ y $B$ por
canales adicionales y la prueba actual no aplica directamente. Tampoco se
prueba bienestar, reducción de volatilidad total ni desaparición de errores
de valoración.

\section{Conclusión}
La señal ruidosa permite derivar las varianzas mediante Bayes.
CARA y normalidad entregan demanda y precio. La sensibilidad esperada
aumenta bajo sentimiento fijo y ausencia de inferencia adicional desde precios.
\label{endbody}
\clearpage
\bibliographystyle{plain}
\bibliography{references}
\clearpage
\appendix
\section{Integral normal para CARA}
Para $W\sim N(M,V)$, $V>0$, completo el cuadrado del exponente:
\begin{equation}
 aw-\frac{(w-M)^2}{2V}
 =-\frac{[w-(M+aV)]^2}{2V}+aM+\frac{a^2V}{2}.
\end{equation}
Extraigo las constantes de la integral:
\begin{align}
 \E[e^{aW}]&=e^{aM+a^2V/2}
 \int_{\mathbb R}\frac{e^{-[w-(M+aV)]^2/(2V)}}{\sqrt{2\pi V}}\,dw\\
 &=e^{aM+a^2V/2}.
\end{align}
La integral vale uno porque es una densidad normal con media $M+aV$ y
varianza $V$. Tomar $a=-g_i$ y el signo negativo de la utilidad da el
resultado. Si $\omega_i=0$, la riqueza es constante y la misma fórmula se
verifica directamente, sin utilizar una densidad de varianza cero.

\section{Diferencia finita}
Sean $d_j(z)=e_j+(1-z)s$ y $B_z=\sum_jn_j/d_j(z)$. Al poner las fracciones
sobre denominador común y simplificar,
\begin{equation}
 \K(y)-\K(x)=
 \frac{(y-x)\left[B_xB_y+
 a\sum_{j=3}^4 n_je_j/(d_j(x)d_j(y))\right]}
 {R(a+B_x)(a+B_y)}>0.
\end{equation}
La regla usada es
$(1-x)n_j/d_j(x)-(1-y)n_j/d_j(y)
=n_je_j(y-x)/[d_j(x)d_j(y)]$. Los denominadores son positivos.
El código verifica la identidad exactamente; Lean prueba la monotonicidad
del coeficiente por una descomposición equivalente.

\section{Correspondencia con Lean y alcance}
La carpeta preliminar conserva las pruebas anteriores y añade
\texttt{SignalCARA.lean}. Sus informes permiten revisar las compilaciones.
La auditoría es parcial: identidades sobre momentos no prueban por sí mismas
Bayes, independencia gaussiana ni una integral de utilidad.
\begin{center}\small
\begin{tabular}{>{\raggedright\arraybackslash}p{1.8cm}>{\raggedright\arraybackslash}p{6.3cm}>{\raggedright\arraybackslash}p{5cm}}
\toprule
Resultado & Declaración / archivo & Alcance\\
\midrule
Prop.~\ref{prop:information}
& \texttt{public\_variance}, \texttt{hidden\_variance},
\texttt{public\_hidden\_covariance}, \texttt{posterior\_square}
en \texttt{SignalCARA.lean}
& Identidades algebraicas probadas. Bayes e independencia probabilística
abiertos en Lean.\\
Prop.~\ref{prop:demand}
& \texttt{cara\_normal\_exponent} en \texttt{SignalCARA.lean};
\texttt{demand\_maximizes}, \texttt{demand\_satisfies\_FOC}
en \texttt{BaselineAudit.lean}
& Exponente, CPO y máximo probados; integral normal abierta en Lean.\\
Prop.~\ref{prop:price}
& \texttt{market\_clearing} en \texttt{BaselineAudit.lean}
& Vaciado algebraico probado; esperanza requiere puente probabilístico.\\
Prop.~\ref{prop:main}
& \texttt{four\_group\_information\_increases},
\texttt{full\_information\_endpoint} en \texttt{Transparency.lean}
& Monotonicidad y extremo probados para el coeficiente algebraico
con parámetros positivos explícitos.\\
\bottomrule
\end{tabular}
\end{center}
No se presenta esta carpeta como una corrida integral de AppliedModelingLib
sobre un commit fijado. Ese procedimiento y el control
\texttt{paper\_contribution.py check ... --fast} quedan pendientes.
El paper final del curso debe reportar el commit formalizado y el resultado
de esa comprobación.

\section{Auditoría de contribuciones de IA}
El registro documenta afirmaciones, veredictos y comprobaciones.
Los diálogos de \texttt{prompts.md} son simulados y están identificados así.
\begin{center}\small
\begin{tabular}{>{\raggedright\arraybackslash}p{4.2cm}>{\raggedright\arraybackslash}p{4.5cm}>{\raggedright\arraybackslash}p{4.3cm}}
\toprule
Afirmación revisada & Veredicto & Comprobación\\
\midrule
Asignar $\mu s$ y $(1-\mu)s$ basta como derivación.
& Incompleto: era una especificación.
& Señal, Bayes, varianzas y covarianza desarrolladas.\\
Media-varianza es numéricamente igual a utilidad esperada CARA.
& Requiere precisión: es un equivalente cierto con igual elección.
& Integral normal y monotonía de la transformación.\\
La transparencia siempre mejora el precio.
& Demasiado amplio: sensibilidad esperada con sentimiento fijo.
& Prueba, barrido racional y Lean algebraico.\\
Información institucional implica filtración del emisor.
& No se desprende del modelo.
& Lectura de secciones 3 y 6.2 y delimitación propia.\\
\bottomrule
\end{tabular}
\end{center}
La guía \texttt{hand/derivaciones-guia.md} organiza las derivaciones para
reproducirlas a mano. Una guía tipografiada no sustituye el apéndice
manuscrito personal exigido por el curso.
\end{document}
